\begin{abstract}
\noindent The records that make social facts are moving out of drawers and registries into systems
that hold them for us: platforms, shared ledgers, and collaborations in which each party keeps its
own log of signed records. Earlier work modeled such collaborations as machines that compose, and
showed, in a form of the CALM theorem, that steps which only add records converge without
coordination while a step that depends on a record being absent needs coordination. This paper asks
what that result means for social ontology. We call the condition on records under which a social
fact holds its grounding condition, and sort grounding conditions by monotonicity. A fact whose
condition asks only that certain records exist can be settled by anyone who holds them, in any
order. A fact that asks for an absence, such as that no objection is open, that a prescription is
the latest, or that a license has not been revoked, cannot be settled by anyone until a seal
declares a set of records complete. A seal is a status function declaration in Searle's sense, and
the parties a question concerns can make it together, as a joint commitment in Gilbert's sense.
Authority is needed where absence matters, and where it is needed it can be shared. We then ask who
may seal, and answer with secondary rules written as co-signed records, frames built in layers by
consensus, tiered thresholds, governing bodies that act through one facade, and appeals that end in
institutions outside the collaboration. Replies to the main objections close the paper.

\medskip\noindent\textbf{Keywords:} social ontology; institutional facts; CALM theorem; joint
commitment; documentality; governance
\end{abstract}
